\FloatBarrier
\part*{Part C: Cross-system material}
\addcontentsline{toc}{part}{Part C: Cross-system material}

\section{Horizon profiles of the baselines and the autocorrelations of Section 5}\label{si:horizon}

Table~\ref{tab:si_horizon} gives the rolling-origin RMSE of last-value persistence and of the training mean at every horizon $h=1,\dots,5$ on the three scored series, computed from the registered data files on the origin sets of the scored ladders (first origin 1990 with an 8-year minimum training window on Specification A; 1965 with 12 years on Specification B; 1948 with 15 years on J-17), together with the lag-$h$ autocorrelation $\rho_h$ of the predictand over the full series and the stationary prediction $\sqrt{2(1-\rho_h)}$ of main-text Remark~2 for the ratio of the two RMSEs. The $h=1$ and $h=5$ values reproduce the archived summaries (main-text Tables~3 and 7) to the printed precision: 98.0/264.7 and 423.6/507.4\,kt on Specification A, 84.4/300.0 and 458.1/521.6\,kt on Specification B, 13.23/21.11 and 16.17/16.80\,ft on J-17. On J-17 the observed ratio tracks the prediction within 0.11 at every horizon, always from below, and crosses one between $h=1$ and $h=2$, where $\rho_h$ falls through one half. On both cod series the observed ratio stays far below the prediction at every horizon, the training mean of a collapsing series being a poor forecast of anything; the stationary formula would place a crossover between four and five years on Specification A and beyond five on Specification B, and neither occurs. The script is \texttt{work/horizon\_profile.py}; its JSON summary is archived beside it.

\begin{table}[!htb]
\centering\small
\caption{Horizon profile of the two baselines on the scored origin sets. RMSE in kt (cod) and ft (J-17); ratio $=$ RMSE(persistence)/RMSE(training mean); $\rho_h$ is the lag-$h$ autocorrelation of the predictand over the full series.}
\label{tab:si_horizon}
\begin{tabular}{@{}lrrrrrrr@{}}
\toprule
Series & $h$ & $n$ & persist & mean & ratio & $\rho_h$ & $\sqrt{2(1-\rho_h)}$ \\
\midrule
J-17 head (ft) & 1 & 75 & 13.23 & 16.17 & 0.818 & 0.629 & 0.861 \\
 & 2 & 74 & 17.69 & 16.58 & 1.067 & 0.348 & 1.142 \\
 & 3 & 73 & 20.59 & 16.88 & 1.220 & 0.137 & 1.314 \\
 & 4 & 72 & 21.50 & 16.93 & 1.270 & 0.061 & 1.371 \\
 & 5 & 71 & 21.11 & 16.80 & 1.256 & 0.073 & 1.361 \\
\midrule
Cod A, SSB (kt) & 1 & 25 & 98.05 & 423.59 & 0.231 & 0.933 & 0.365 \\
 & 2 & 24 & 184.32 & 460.54 & 0.400 & 0.825 & 0.592 \\
 & 3 & 23 & 237.97 & 488.07 & 0.488 & 0.687 & 0.792 \\
 & 4 & 22 & 256.74 & 500.18 & 0.513 & 0.553 & 0.945 \\
 & 5 & 21 & 264.72 & 507.39 & 0.522 & 0.406 & 1.090 \\
\midrule
Cod B, SSB (kt) & 1 & 59 & 84.43 & 458.12 & 0.184 & 0.971 & 0.241 \\
 & 2 & 58 & 152.44 & 474.51 & 0.321 & 0.913 & 0.417 \\
 & 3 & 57 & 208.89 & 489.99 & 0.426 & 0.838 & 0.569 \\
 & 4 & 56 & 258.77 & 505.80 & 0.512 & 0.757 & 0.697 \\
 & 5 & 55 & 299.98 & 521.56 & 0.575 & 0.678 & 0.802 \\
\bottomrule
\end{tabular}

\end{table}

Table~\ref{tab:si_autocorr} gives the lag-one autocorrelations, means and standard deviations used in main-text Section~5.1. The J-17 flux mechanism is $\hat\beta R_t+\hat\gamma P_t$ with the full-sample least-squares coefficients of the one-pool map, $(\hat\alpha,\hat\beta,\hat\gamma,\hat\delta)=(177.6,\,0.0173,\,-0.0260,\,-0.273)$; the apparent net production of the cod series is $S_{t+1}-S_t+C_t$ with the catch series of the respective pass. The stylised prediction of Proposition~2(ii) for J-17 uses $\sigma^2=84.6$\,ft$^2$ and $\rho=0.078$, giving $\sigma^2(1-2\rho)=71.4$\,ft$^2$, a predicted one-year M2 RMSE of $\sqrt{13.23^2+71.4}=15.70$\,ft against the archived 14.70\,ft, and a predicted M2m-over-M2 mean-square advantage of 71.4\,ft$^2$ against the archived $14.70^2-12.28^2=65.1$\,ft$^2$. Further correlations used in the main text: $r(\Delta H_t,R_t)=0.742$, $r(\Delta H_t,R_{t-1})=0.056$, $r(R_t,R_{t-1})=0.172$ as a Pearson correlation of successive pairs (0.170 as the autocovariance-based estimator in the table), and $r(H_t,H_{t-1})=0.644$ as a Pearson correlation of pairs (0.629 autocovariance-based; the single-system aquifer study quoted 0.64).

\begin{table}[!htb]
\centering\small
\caption{Lag-one autocorrelations (autocovariance estimator), means and standard deviations of the series used in main-text Section~5.1. Units: ft and $10^3$ acre-ft for J-17; kt for cod.}
\label{tab:si_autocorr}
\begin{tabular}{@{}lrrr@{}}
\toprule
Series & $\rho_1$ & mean & s.d. \\
\midrule
J-17: annual recharge $R_t$ & 0.170 & 682.83 & 515.44 \\
J-17: annual pumpage $P_t$ & 0.870 & 318.82 & 114.41 \\
J-17: flux mechanism $\hat\beta R_t+\hat\gamma P_t$ & 0.078 & 3.52 & 9.20 \\
J-17: head increment $\Delta H_t$ & -0.115 & -0.39 & 12.34 \\
J-17: head $H_t$ & 0.629 & 664.42 & 14.70 \\
Cod A: apparent net production, coarse regime & 0.753 & 57.74 & 112.71 \\
Cod A: apparent net production, annual landings & 0.714 & 53.00 & 124.86 \\
Cod B: apparent net production, landings & 0.908 & 184.73 & 212.65 \\
Cod A: SSB $S_t$ & 0.933 & 300.51 & 356.51 \\
Cod B: SSB $S_t$ & 0.971 & 512.76 & 423.25 \\
\bottomrule
\end{tabular}

\end{table}

\section{Verification record for this edition}\label{si:verify}

The verification script \texttt{paper11\_forecasting\_\allowbreak{}baselines\_v67\_\allowbreak{}verification.py} (standard library only) reads the registered data files and archived result files of both systems together with the manuscript sources (main text and this document) and asserts, in order: (a) the cod outcome-year, vintage and threshold cross-checks carried over from the single-system edition (S10--S12), including the presence of every pinned value in the sources; (b) every cell of main-text Tables~7--9 against \texttt{wave\_e\_edwards/results/rolling\_summary.csv}, \texttt{pass2\_H\_summary.csv}, \texttt{pass2\_R\_summary.csv} and \texttt{e3\_audit\_uncertainty.json}, the fixed-window cells of S17 against \texttt{fixed\_window\_scores.csv}, the post-2007 record of S19 against \texttt{rolling\_modern\_2007.csv} and the Comal ranking of S20 against \texttt{fibre\_comal\_summary.csv}; (c) the baseline horizon profile of S23 recomputed from the registered series, its agreement with the archived $h=1$ and $h=5$ summaries of both systems, and the autocorrelations quoted in Section~5; (d) the arithmetic of Proposition~2 on the J-17 numbers (the sign of $\sigma^2(1-2\rho)-\mu^2$, the predicted ordering M2m $\le$ persist $<$ M2, and the agreement of the predicted M2m-over-M2 advantage with the archived one to within 10\%); and (e) textual invariants (no claim of prospective registration; the two protocols stated with their disclosed completions; the retraction of the ``weakly autocorrelated surplus production'' reading accompanied by the computed value). The run for this edition reports 139 checks passing and none failing; its output is archived beside the script.

\end{document}
